\section{ADR-010: Practice analysis and time-stretch}
\textbf{Status:} accepted and implemented (M14).

\textbf{Analysis.} A pipeline step \texttt{analyse} runs after decoding, before
separation. Its cache key is the decoded mix's SHA-256 plus the algorithm
version (\texttt{analysis-v1}), and the result is a row of the
\texttt{analyses} table (migration 6), so the same mix is never analysed twice;
\texttt{GET /items/\{id\}/analysis} computes it on demand for songs separated
before M14. Everything uses librosa (ISC) and classic signal processing, so no
model weights are needed:
\begin{itemize}
  \item \emph{tuning}: \texttt{librosa.estimate\_tuning}, in cents from A440;
  \item \emph{key}: mean constant-Q chroma of the harmonic part (HPSS keeps drum
  noise out), corrected for the tuning, correlated with the 24 rotations of the
  Krumhansl--Kessler and Temperley key profiles, averaged; the confidence is
  the margin over the runner-up;
  \item \emph{tempo and beats}: \texttt{librosa.beat.beat\_track}; the BPM is
  the median inter-beat interval;
  \item \emph{downbeats}: librosa has none, so 4/4 is assumed and the bar phase
  is the beat offset (0--3) with the most low-frequency onset energy.
\end{itemize}
Beat This! (MIT code and weights) tracks beats and downbeats better, but needs
PyTorch and a weight download; it stays a possible upgrade behind the same
API. Named section markers are stored per mix in the \texttt{markers} table.

\textbf{Real-time pitch and tempo.} The preview keeps one
\texttt{AudioBufferSourceNode} per track (ADR-007). A worklet cannot pull those
faster, so the tempo is the sources' \texttt{playbackRate} (a resample, which
also raises the pitch by $12\log_2 r$), and one Signalsmith Stretch node (MIT,
official WebAssembly build) on the summed bus transposes by
$s + c/100 - 12\log_2 r$ semitones. One instance serves every stem, so the cost
does not grow with the number of stems: an offline render of six stems through
the chain took 2.8\,\% of real time on one core of the i5-11400H in Chromium
(0.4\,\% without the stretch). The node is only in the chain when the shift is
not zero. Its module is served from the app's origin and the CSP gains only
\texttt{'wasm-unsafe-eval'} (no \texttt{blob:} scripts). Rubber Band (GPL) and
SoundTouch (LGPL) were rejected for their licences.

\textbf{Exports.} \texttt{POST /items/\{id\}/render} takes \texttt{tempo},
\texttt{semitones}, \texttt{cents}, \texttt{click\_gain} and
\texttt{with\_stems}. The engine applies the same two steps in the same order:
resample by the tempo (soxr through librosa), then Signalsmith Stretch through
\texttt{python-stretch} (MIT) with the same tonality limit (8~kHz), so the
export matches the preview. The click stem is synthesised from the stored beat
grid with the same recipe in the UI and the engine (short 1~kHz blips, 1.5~kHz
and louder on downbeats).

\textbf{Consequences.} The stretch adds about 60~ms of latency to the preview
when active, and the position display does not compensate for it. The
resamplers differ (the browser's against soxr), which is inaudible next to the
stretch. Songs in other metres get 4/4 bars until a downbeat tracker is added.
